\section{Timing and transfer}
We isolated the arithmetic cost of voting by giving two implementations the
same columns, seeds, pruning rules and tie order. The reference caches direct
Cholesky fits and reuses the full Gram matrix. At seed 42 the two methods
returned identical indices, rankings and first negative witnesses on all 14
sources tested. On six previously studied sources with 72--100 columns,
sharing the base fit gave a geometric mean $8.73$-fold ranking speedup.
Seven of the eight newer sources have at most nine columns. Here direct fits
were faster on seven sources and the direct-to-batched time ratio averaged
$.94$ geometrically. Matrix updates for fast regularized least-squares
selection have a substantial prior literature~\cite{pahikkala2010rls,pahikkala2016rlscore}.
The identical-game comparison measures their value for this sampling schedule.

The six wider classification sources supplied 16 archived folds. We preserved
their training, voting, width and test rows and used at most the first 100 columns
in dataset order. The primary rule retained 32.8\% of the available columns
geometrically. Mean AUC and accuracy changed by $+.0016$ and $+.0080$ relative
to full ridge. Compared with unit-step RFE, MinPI took 44.5\% of the selection
time and retained $1.95$ times as many columns. Against coefficient ranking it
took $1.21$ times as long and retained $1.08$ times as many columns. These
comparisons revisit known sources. In particular, truncating CNAE to its first
100 columns left full-model accuracy at .302.

\paragraph{The original six benchmark sources.}
The prior experiments used a GPU histogram booster with 100 depth-four trees,
ranking columns by mean split gain~\cite{chen2016xgboost}. On the original
prepared splits, all 30 gain orders matched the archive, including ties.
Both rankers use the original training rows; MinPI also reads the reserved
voting outcomes. Gain uses an RTX 5070 Laptop GPU and ridge one host thread.
Table~\ref{tab:gain-receiver} pairs each ranker with a common receiver for
width validation and final prediction. The game used for voting always refits
ridge. On the five classification sources the 95\% source-resampling intervals
for the MinPI-to-gain width ratio were $[.72,1.20]$ with ridge and
$[1.18,1.77]$ with boosted trees. On 25 of the 30 ridge splits, coefficient
ranking alone selected the same columns as MinPI.

\begin{table}[h]
\centering\small
\caption{Sampled MinPI relative to original CUDA gain on the same candidate
widths. Time ratios cover ranking and width validation. Width and time are
paired geometric ratios. Predictive differences subtract the gain result.
The five classification sources receive equal weight; regression is Abalone
alone. Each source has five splits. Ridge clocks use three timed repeats,
XGBoost clocks one.}
\label{tab:gain-receiver}
\begin{tabular}{@{}llrrrr@{}}
\toprule
Receiver & Domain & Time ratio & Width ratio & $\Delta$AUC or $\Delta R^2$ & $\Delta$Acc.\ or $\Delta$MAE\\
\midrule
Ridge & Classification & .272 & .906 & $+.00001$ & $-.00126$\\
XGBoost & Classification & .896 & 1.467 & $-.00357$ & $-.00511$\\
Ridge & Abalone & .097 & .977 & $+.00742$ & $-.01404$\\
XGBoost & Abalone & .920 & 1.516 & $+.00286$ & $-.00650$\\
\bottomrule
\end{tabular}
\end{table}

On ISOLET both rankings returned 75 columns on average with XGBoost.
MinPI and gain gave .829 and .827 accuracy and took 26.75 and 30.28 seconds
to select. On Madelon all 500 feature--split indices were sampled zeros.
MinPI returned 40.2 columns against gain's 20.4 and its accuracy fell from
.758 to .730. The previously deployed gain pipeline, which examined every
prefix with XGBoost, gave .846 AUC there; replacing both the ranking and
the predictor with ridge MinPI gave .624. Across the five classification
sources that replacement cost .058 AUC and 7.44 accuracy points.
Using XGBoost for both validation and deployment reduced these gaps to
.009 AUC and .89 points. The receiver follow-up reused the historical
splits. Excluding test rows whose prepared predictor vector occurred in
development gave similar paired differences.
